\subsection{Properties of Inequality}
In addition to equality, we can describe other relationships between expressions.
\begin{center}
\begin{tabularx}{\textwidth-1pt}{|l|>{\hsize=0.375in\centering\arraybackslash}X|>{\hsize=0.375in\centering\arraybackslash}X|X|}
\hline
\textbf{Name}        &\multicolumn{2}{|c|}{\textbf{Symbol}}&\textbf{When is it true?}\\\hline
Less Than            &$<$   &&The expression on the left represents a lesser number than the one on the right.\\\hline
Less Than or Equal   &$\leq$&&The expression on the left represents a lesser or the same number.\\\hline
Greater Than         &$>$   &&The expression on the left represents a greater number than the one on the right.\\\hline
Greater Than or Equal&$\geq$&&The expression on the left represents a greater or the same number.\\\hline
Not Equal            &$\neq$&&The two expressions are not equal.\\\hline
\end{tabularx}
\end{center}
%We say the \term{lesser} number is the one further to the left on the number line, and the \term{greater} number is the one further to the right on the number line.
%\begin{center}
%\begin{tikzpicture}
%\draw[thick] (-10,0) -- (10,0);
%\draw[thick] (-9.8,0.2) -- (-10,0) -- (-9.8,-0.2);
%\draw[thick] (9.8,0.2) -- (10,0) -- (9.8,-0.2);
%\foreach \x in {0,1,...,9} {
	%\draw (\x,0.2) -- (\x,-0.2);
	%\draw (\x,-0.8) node[anchor=base]{$\x$};
%}
%\foreach \x in {1,...,9} {
	%\draw (-\x,0.2) -- (-\x,-0.2);
	%\draw (-\x,-0.8) node[anchor=base,xshift=-0.75ex]{$-\x$};
%}
%\end{tikzpicture}
%\end{center}
%\begin{Example}
%Fill in each blank below with an appropriate inequality symbol.
%\begin{colparts}{3}
%\foreach \x/\a/\b in {0/2/5,1/-3/-7,2/-1/-1} {
	%\regtext{\x}{0}{$\a\;\makeblanksmall\;\b$};
%}
%\end{colparts}
%\end{Example}
\begin{definition}[Inequality]
An \term{inequality} is a statement describing a relationship between expressions using an inequality symbol.

A \term{solution} to an inequality is a value(s) of the variable(s) in the inequality that makes the inequality true.
The \term{solution set} of an inequality is the set of all solutions of an inequality.
Two inequality are \term{equivalent} if they have the same solution set.
To \term{solve} an inequality means to find its solution set.
\end{definition}
Like equations, we have properties which allow us to get equivalent inequalities. The properties are shown here for $<$, but they're paralleled for other inequality symbols.
\begin{colparts}{2}
\draw[use as bounding box] (-1,0) rectangle (1,-12);
\draw (-1,-6.5) -- (1,-6.5);
\draw (0,0) -- (0,-12);
\foreach \x/\lbl/\lblb/\eqn in {-0.5/Positive/c>0/c(X)<c(Y),0.5/Negative/c<0/c(X)>c(Y)} {
	\regtext{\x}{-1}{\textbf{Multiplicative Property of}};
	\regtext{\x}{-2}{\textbf{Inequality (\lbl)}};
	\regtextleftblock{\x}{-3}{For any expressions $X$ and $Y$ and any real number $\lblb$, the inequalities}{0.5\textwidth-0.25in};
	\regtext{\x}{-5}{$X<Y$\quad and\quad$\eqn$};
	\regtextleftblock{\x}{-6}{are equivalent.}{0.5\textwidth-0.25in};
}
\foreach \x/\lbl/\lblb/\eqn in {-0.5/Additive/ and any real number $c$/X+c<Y+c,0.5/Antisymmetric/{}/Y>X} {
	\regtext{\x}{-7.5}{\textbf{{\lbl} Property of Inequality}};
	\regtextleftblock{\x}{-8.5}{For any expressions $X$ and $Y$\lblb, the inequalities}{0.5\textwidth-0.25in};
	\regtext{\x}{-10.5}{$X<Y$\quad and\quad$\eqn$};
	\regtextleftblock{\x}{-11.5}{are equivalent.}{0.5\textwidth-0.25in};
}
\end{colparts}
